\apendice{Manual de especificación de diseño}

\section{Planos}
Este proyecto es una aplicación web de software, por lo que no dispone de planos físicos de hardware ni esquemas de distribución de dispositivos.


\section{Diseño arquitectónico}
La aplicación sigue el patrón arquitectónico \textbf{MVC (Modelo-Vista-Controlador)} implementado sobre el framework Flask en Python.

\begin{description}
 	\item[\textbf{Modelo}] Gestiona la persistencia de datos mediante SQLAlchemy como ORM. El modelo de datos está compuesto por la tabla central \texttt{saac\_sistema}, las tablas maestras de referencia (\texttt{tipo\_entrada}, \texttt{idioma}, \texttt{plataforma}, 					\texttt{entorno\_uso}, \texttt{metodo\_comunicacion}), las tablas de relación muchos a muchos (\texttt{sistema\_entrada}, \texttt{sistema\_idioma}, \texttt{sistema\_plataforma}, \texttt{sistema\_metodo}, \texttt{sistema\_entorno}, \texttt{sistema\_dependencia}), la tabla de requisitos 		funcionales (\texttt{sistema\_requisito\_funcional}) y la tabla de historial anónimo (\texttt{historial\_recomendacion})
 	\item[\textbf{Vista}] Las plantillas HTML son renderizadas por el motor Jinja2. Las vistas principales son el formulario de cuestionario, la pantalla de resultados segmentada en tres bloques (sistemas finales, sistemas preventivos y bancos de voz), el formulario de feedback y el panel 		de administración.
 	\item[\textbf{Controlador}] La lógica de negocio reside en las rutas Flask: \texttt{/recomendar} procesa el cuestionario y ejecuta el algoritmo de filtrado clínico, \texttt{/guardar-feedback} vincula la retroalimentación del usuario con su registro en base de datos, y \texttt{/				admin\_historial} expone el panel de consulta para el investigador.
\end{description}

\imagen{../img/diagrama_arquitectura.png}{Diagrama de arquitectura de la aplicación.}

\subsection{Diseño de vistas}
La aplicación cuenta con las siguientes vistas, organizadas según su naturaleza.

\subsubsection{Páginas informativas estáticas}
Accesibles desde el menú de navegación principal, proporcionan contexto clínico y legal al usuario antes de interactuar con el cuestionario.

\begin{description}
		\item - Página de inicio: Presenta la naturaleza y propósito de la herramienta, los criterios utilizados para su desarrollo y las instituciones colaboradoras.
		\item- Información sobre ELA: Explica qué es la Esclerosis Lateral Amiotrófica, sus tipos según el origen, sus fases de progresión y consideraciones sobre la preservación cognitiva.
		\item- Información sobre SAAC: Describe los Sistemas Aumentativos y Alternativos de Comunicación, su clasificación según tecnología, la importancia de los sistemas escalables y aspectos sobre financiación.
		\item- Base legal: Detalla la finalidad del tratamiento de datos, las bases jurídicas, la política de conservación, los destinatarios de la información y los derechos del usuario en cumplimiento del RGPD.
\end{description}

\subsubsection{Páginas funcionales:}
Constituyen el flujo principal de interacción con el sistema.

\begin{description}
		\item- Cuestionario: Formulario web que recoge el perfil clínico del usuario mediante menús desplegables y botones de opción, incluyendo la casilla de consentimiento RGPD obligatorio.
		\item- Resultados: Renderiza dinámicamente la propuesta personalizada segmentada en tres bloques: sistemas finales, sistemas de uso preventivo y bancos de voz.
		\item- Feedback: Formulario opcional para que el usuario evalúe la propuesta recibida.
		\item- Agradecimiento: Pantalla de confirmación tras el envío del feedback.
		\item- Panel de administración: Vista exclusiva para el investigador que lista cronológicamente todos los registros almacenados con consentimiento, diferenciando casos reales de pruebas.
\end{description}

\subsection{Diseño de la base de datos}
La base de datos relacional está compuesta por las siguientes tablas, organizadas en tres grupos funcionales.

\subsubsection{Tablas maestras}
Almacenan los valores de referencia del sistema de filtrado clínico.

\begin{description}
	\item[\texttt{tipo\_entrada}] Métodos de entrada disponibles: manos, ojos, cabeza, voz y pulsador.
	\item[\texttt{idioma}] Idiomas soportados: español, gallego, catalán y euskera.
	\item[\texttt{plataforma}] Plataformas tecnológicas: Windows, iOS, Android, Web y panel físico.
	\item[\texttt{entorno\_uso}] Contextos de uso: domicilio, exterior y mixto.
	\item[\texttt{metodo\_comunicacion}] Métodos de comunicación: alfabeto y pictogramas.
\end{description}

\subsubsection{Tablas principales de sistemas}
\begin{description}
	\item[\texttt{saac\_sistema}] Tabla central del catálogo. Almacena todos los sistemas, hardware y servicios recomendables, con sus atributos: nombre, descripción, si requiere interlocutor, nivel de fatiga física, portabilidad, compatibilidad con anclaje, enlace externo y categoría 			(sistema, hardware o servicio).
	\item[\texttt{sistema\_requisito\_funcional}] Define los niveles mínimos funcionales que debe tener el usuario para poder utilizar cada sistema: nivel visual, auditivo, tecnológico y de habla, en una escala de 0 a 3.
\end{description}

\subsubsection{Tablas de relación (muchos a muchos)}
Vinculan cada sistema con sus características clínicas y técnicas.

\begin{description}
	\item[\texttt{sistema\_entrada}] Relaciona cada sistema con los métodos de entrada que admite.
	\item[\texttt{sistema\_idioma}] Relaciona cada sistema con los idiomas en los que está disponible.
	\item[\texttt{sistema\_plataforma}] Relaciona cada sistema con las plataformas en las que funciona.
	\item[\texttt{sistema\_metodo}] Relaciona cada sistema con el método de comunicación que emplea.
	\item[\texttt{sistema\_entorno}] Relaciona cada sistema con los entornos de uso en los que es apto.
	\item[\texttt{sistema\_dependencia}] Define qué periféricos o hardware requiere cada sistema para funcionar, indicando además si dicho hardware es obligatorio u opcional.
\end{description}

\subsubsection{Tabla de historial}
\begin{description}
	 \item[\texttt{historial\_recomendacion}] Almacena de forma anónima los registros generados con consentimiento de investigación. Sus campos principales son: identificador único autogenerado (\texttt{id}), fecha del registro, perfil clínico del usuario en formato JSONB (\texttt{input\_usuario}), sistemas recomendados, retroalimentación estructurada en JSON (\texttt{feedback}) e indicador booleano (\texttt{es\_real}) que distingue casos clínicos reales de pruebas.
\end{description}

\subsection{Flujo de datos}
El flujo principal de la aplicación sigue estos pasos:

\begin{enumerate}
	\item El usuario navega por las páginas informativas (inicio, ELA, SAAC y base legal) para conocer la herramienta.
	\item El usuario completa el cuestionario y acepta el consentimiento RGPD.
	\item El servidor recibe la petición POST en \texttt{/recomendar} y ejecuta el algoritmo de filtrado clínico.
	\item Si el usuario otorgó consentimiento de investigación, el perfil se almacena de forma anónima en la base de datos.
	\item Se renderiza la pantalla de resultados con los sistemas SAAC recomendados.
	\item Opcionalmente, el usuario envía feedback, que queda vinculado a su registro mediante el \texttt{historial\_id}.
	\item El investigador puede consultar todos los registros almacenados accediendo a \texttt{/admin\_historial}.
\end{enumerate}